% ---------------------------------------------------------------------------
% What the master program contains, and how a recovered design is scored.
% Placed immediately after the capability matrix: the matrix says what prior
% work cannot express; this says what our program does express, and how we
% evaluate it now that the outer proposal loop is gone.
% ---------------------------------------------------------------------------
\section{What the program contains and how it is scored}
\label{sec:program_contents}

\subsection{Evaluating a design without a proposal loop}
\label{sec:eval_without_vizier}

Vizier proposed hardware configurations; the performance model and execution
evaluator determined their scores. Removing the proposer does not remove that
evaluator. The formulation minimizes modeled latency under hardware
constraints, including area and power budgets, with compute tiles, bandwidth,
SRAM and vector lanes as shared LP variables; discrete physical choices still
require a supported hardware menu and a recovery step. For each recovered chip
and schedule we report the quantities in \cref{tab:eval_metrics}.

\begin{table}[t]
\centering
\small
\caption{Quantities evaluated for each recovered chip and schedule.}
\label{tab:eval_metrics}
\begin{tabularx}{\linewidth}{@{}lX@{}}
\toprule
\textbf{Metric} & \textbf{Evaluation} \\
\midrule
Prefill     & Prompt execution latency and time to first output \\
Decode      & Remaining generation latency and time per output token \\
Total       & Actual full-task makespan \\
Throughput  & Completed output tokens divided by elapsed time \\
Perf/TDP    & Throughput divided by the chip's TDP \\
Speedup     & Matched baseline latency divided by candidate latency \\
\bottomrule
\end{tabularx}
\end{table}

For a fixed amount of work these compose exactly,
\begin{equation}
\frac{(\mathrm{Perf}/\mathrm{TDP})_{\mathrm{design}}}
     {(\mathrm{Perf}/\mathrm{TDP})_{\mathrm{baseline}}}
=
\frac{T_{\mathrm{baseline}}}{T_{\mathrm{design}}}\,
\frac{\mathrm{TDP}_{\mathrm{baseline}}}{\mathrm{TDP}_{\mathrm{design}}},
\label{eq:perf_tdp_decomp}
\end{equation}
so a latency objective does not by itself maximize the ratio. We obtain a
performance/power tradeoff by solving across power budgets or hardware
branches and comparing the recovered designs, rather than by optimizing
\cref{eq:perf_tdp_decomp} directly.

Three limits apply to every number reported this way. These are analytical
execution results: physical timing and power require calibration against
implementations, simulation or measurement. TDP is not average consumed power,
so energy per inference needs a separate power calculation. And the accepted
prefill result remains our reference point --- new total-inference claims
require a matched workload and numerical contract, and decode's share must come
out of the evaluated schedule rather than be imposed as a target fraction.

\subsection{Features represented in the master}
\label{sec:master_features}

One implementation distinction matters for reading the results that follow. The
master architecture carries several branches and extensions, and the serial-policy
solver experiment of \cref{sec:experiments} does not activate all of them
together. \cref{tab:master_features} is an inventory of the available
formulations, not a claim that every row is exercised in every reported solve.

\begin{table}[t]
\centering
\scriptsize
\setlength{\tabcolsep}{4pt}
\caption{Formulations represented in the master program. Rows marked
\emph{extension} are implemented and tested but not activated in the
serial-policy experiment.}
\label{tab:master_features}
\begin{tabularx}{\linewidth}{@{}lX@{}}
\toprule
\textbf{Feature} & \textbf{What it represents} \\
\midrule
Original paper V1 & Supplied complete configuration, placement, retention and rematerialization programs, preserving their original costs \\
Independent mapping and tiling & Separate implementation choices for logical invocations, including different decode positions \\
Buffering/dataflow modes & Choices supplied by each qualified implementation menu \\
Batching & Compatible operations may share a batch implementation while retaining their individual identities \\
Fusion & Supported fusion groups and fused kernel atoms, including SwiGLU variants \\
Temporal scheduling & Start, work, finish and active state over event labels \\
Dependency enforcement & Producer completion, operand readiness and control dependencies \\
Routing causality & Expert work becomes available only after the relevant router dependency \\
Time-dependent placement & Host, DRAM and SRAM presence throughout execution \\
Retention and eviction & Keep useful data between consumers or release its storage \\
Explicit transfers & Charged host$\leftrightarrow$DRAM and DRAM$\leftrightarrow$SRAM copies with source and destination obligations \\
Weight residency/streaming & Retained weights versus supported streamed-weight implementations \\
Streamed inputs/outputs & Qualified tile staging instead of requiring every operand/output to be fully resident \\
KV across timesteps & KV production, persistence, subsequent reads, growth and declared terminal backing \\
KV spill/restore & Host backing, reservations, transfer completion and readiness \emph{(extension)} \\
Rematerialization & Recompute declared deterministic values, charging inputs, execution and temporary storage \\
Typed representations/conversions & Distinct stored/runtime representations, conversion work, destination allocation and readiness \emph{(extension)} \\
Capacity accounting & Live tensors, incoming copies, workspace, output reservations and conversion temporaries \\
Shared resource accounting & Compute, vector, local-memory, bandwidth and host-link service \\
Shared hardware co-design & One common hardware allocation for the execution \\
Hardware constraints & Bounds plus area, peak-power and power-density proxies \\
Linear relaxations & Reciprocal envelopes and product relaxations connecting hardware rates to service time \\
\bottomrule
\end{tabularx}
\end{table}

The program-sheet inventory and the conversion and spill/restore extension
records document each row; \cref{tab:master_rows} gives the row families that
realize them.
